\subsection{A scalar auxiliary place and a proper ordinary determinant quotient}
\label{sec:cubic-scalar-auxiliary-support}

This section works at \(p=269\) and \(K=\mathbb Q(\sqrt{-3})\).
It retains a completed instance of AC-G78--AC-G83 and AC-G87,
with its actual final CM host \(H\), exact old inertia types,
ordered coefficient-prime flags, and separated residual characters.
The original Dirichlet modulus is denoted by \(N_\psi\).
It is not the tame level \(N_f\) of the ordinary modular form:
the latter can be divisible by \(3\).
The finite character \(\psi\) and integer \(e\geq1\) satisfy the
literal original \(\mathrm{GD}\) conditions
\[
 p\nmid\operatorname{ord}(\psi),\qquad
 \operatorname{comp}_p(\psi)^2=1,\qquad
 \psi(-1)=(-1)^e,\qquad (3,N_\psi)=1,\qquad
 268\nmid2e.                                      \tag{G89.1}
\]
Write \(\overline\epsilon\) for the residual cyclotomic character,
\(\overline\gamma=\overline\psi\,\overline\epsilon^e\), and
\(\overline\chi\) for the fixed cubic character on \(G_K\).
The source normalization of the residual parameter on \(G_H\) is
\[
 \overline r=
 \overline\epsilon^{-1}
 \bigl(\overline\chi^{\,j}\overline\gamma
       \oplus1\oplus
       \overline\chi^{\,j}\overline\gamma^{-1}\bigr),
 \qquad j\in\{1,2\}.                              \tag{G89.2}
\]
All residual characters are viewed in one sufficiently large finite
coefficient field.

\begin{lemma}[AC-G89-v0: a source-compatible scalar Frobenius place]
\label{lem:G89-v0}
For either actual host \(H\) of AC-G83, and after excluding any
prescribed finite set of rational primes, there is a rational prime
\(v_0^{\mathbb Q}\) which splits completely in \(H\), is unramified
in all the fixed residual fields, and satisfies
\[
 \overline\epsilon(\operatorname{Frob}_{v_0^{\mathbb Q}})=-1,
 \qquad
 \overline\gamma(\operatorname{Frob}_{v_0^{\mathbb Q}})=1,
 \qquad
 \overline\chi(\operatorname{Frob}_{v_0^{\mathbb Q}})=1.
                                                        \tag{G89.3}
\]
Here the Frobenius is restricted to \(G_K\) for the last equality.
Either arithmetic or geometric Frobenius gives these same values.

Choose a place \(v_0\) of \(H^+\) above this prime.
Then \(\operatorname{Norm}(v_0)=v_0^{\mathbb Q}\), the place is split
in \(H/H^+\), and the residual datum (G89.2), on the source
\(\mathrm{GFrame}\) for this field and place, satisfies the literal
\(\mathrm{Resid.V0ok}\).
\end{lemma}

\begin{proof}
Put \(N_\psi=p^a m\), with \((m,p)=1\).
The source definition of \(\operatorname{proj}(p,N_\psi)\)
is the Chinese-remainder endomorphism which is the identity on
\((\mathbb Z/p^a)^\times\) and is \(1\) on
\((\mathbb Z/m)^\times\).  Thus
\[
 \tau=\operatorname{comp}_p(\psi),\qquad
 \beta=\psi\tau^{-1}                               \tag{G89.4}
\]
are exactly the two conductor components.
The character \(\beta\) factors through \((\mathbb Z/m)^\times\).
If \(a=0\), then \(\tau=1\).
Otherwise the first and second conditions in (G89.1) imply that
\(\tau\) is trivial or quadratic and factors through
\((\mathbb Z/p)^\times\): the kernel of reduction from
\((\mathbb Z/p^a)^\times\) is a \(p\)-group.
Since \(p\equiv1\pmod4\), one has \(\tau(-1)=1\).
Consequently
\[
 (m,3p)=1,\qquad
 \beta(-1)=(-1)^e,\qquad
 p\nmid\operatorname{ord}(\beta).                  \tag{G89.5}
\]
The first assertion uses the original conductor condition
\((3,N_\psi)=1\), not any condition on \(N_f\).

Let \(B_\beta/\mathbb Q\) be the abelian field cut out by
\(\overline\beta\), and let
\[
 E_{269}=\mathbb Q(\sqrt{-3},\sqrt[3]{269}),\qquad
 \Omega=B_\beta\,\mathbb Q(\mu_p)\,E_{269}.          \tag{G89.6}
\]
The extension \(\Omega/\mathbb Q\) has degree prime to \(p\).
The first field is ramified only at primes dividing \(m\);
the other two are ramified only at \(3p\).
Their intersection is therefore unramified at every finite prime
of \(\mathbb Q\), and is \(\mathbb Q\).
Here the usual discriminant bound excludes a nontrivial everywhere
unramified extension of \(\mathbb Q\).
Also \(E_{269}\cap\mathbb Q(\mu_p)=\mathbb Q\):
an abelian subextension of the \(S_3\)-extension \(E_{269}\) is
contained in \(K\), whereas \(K\) ramifies at \(3\).
It follows that there is an element \(s\in\operatorname{Gal}(\Omega/\mathbb Q)\)
whose restrictions are complex conjugation on \(B_\beta\),
the element \(-1\) on \(\mu_p\), and the identity on \(E_{269}\).
It has
\[
 s|_K=1,\quad
 \overline\psi(s)=\overline\beta(-1)=(-1)^e,\quad
 \overline\gamma(s)=1,\quad \overline\chi(s)=1.     \tag{G89.7}
\]
In the equality for \(\overline\psi\), its \(p\)-component has
value \(\tau(-1)=1\).

We next check compatibility with the actual host, rather than
postulating a new linearly disjoint field.
Write \(P\) for the cyclic \(p\)-rotation subgroup in AC-G83.
In the odd branch \(H=FT\), with \(F=KK'\) and
\(T/\mathbb Q\) dihedral with quadratic field \(E\).
Then \(H^P=FE\).
In the positive-sign branch \(H=LT\), with
\(L=KD_0D_1\), \(E=\mathbb Q(\sqrt{r_0r_1})\subset L\),
and \(H^P=L\).
In each case \(P\) is contained in the commutator subgroup:
a quadratic lift acts by inversion, and
\([c,a]=a^{-2}\) generates the same odd-order cyclic group as \(a\).
Every quotient of \(\operatorname{Gal}(H/\mathbb Q)\) of order
prime to \(p\) therefore factors through the displayed
multiquadratic field.

The primes used for \(K'\), \(D_0,D_1\), and the odd-branch \(E\)
were chosen fresh for all the original residual data.
Each quadratic subfield of \(H^P\) other than \(K\) ramifies
at at least one of these fresh primes.
More explicitly, in the odd case any quadratic character other
than the one of \(K\) involves \(K'\) or \(E\), and hence ramifies
at \(r_1\) or \(r_2\); in the other case it involves \(D_0\) or
\(D_1\), and hence ramifies at \(r_0\) or \(r_1\).
The old field \(\Omega\) is unramified there.
Since \([\Omega:\mathbb Q]\) is prime to \(p\), this proves
\[
                         H\cap\Omega=K.           \tag{G89.8}
\]
Thus the identity on \(H\) and \(s\) on \(\Omega\) define one
automorphism of \(H\Omega\).
Chebotarev supplies infinitely many primes with its conjugacy
class, so we may exclude the finite set in the statement.
The required values are conjugacy invariant.  This proves
(G89.3), complete splitting in \(H\), and the norm assertion.

At the selected decomposition group, every character in (G89.2)
is unramified and takes the value \(-1\) at Frobenius.
They are therefore equal on the whole decomposition group.
The same Frobenius has cyclotomic value \(-1\ne1\).
These are precisely the three clauses of
\(\mathrm{Resid.V0ok}\): trivial residual inertia, equality of all
residual characters on \(D_0\), and a nontrivial cyclotomic
value in \(D_0\).
\end{proof}

\begin{lemma}[AC-G89-N4: scalar residual characteristic polynomial suffices]
\label{lem:G89-N4}
Let \(r\) be any source residual datum of positive rank with
\(\mathrm{Resid.V0ok}\).  Then every literal source
\(\mathrm{Point}(r)\) satisfies \(\mathrm{N4}\).
No matrix-scalar residual lifting hypothesis is required.
\end{lemma}

\begin{proof}
Let \(y\) be a source point over \(A=k'[[t]]\), and let
\(g\in D_0\) have cyclotomic value \(c\ne1\) in the residue field.
The common value of the residual characters at \(g\) is some
\(\mu\in k'^\times\).
The field \(k'\) may be larger than the original residue field;
both inequalities remain valid after this extension.
Put \(M=\rho_y(g)\in\operatorname{GL}_n(A)\).
The literal \(\mathrm{Point.red}\) equality of determinants gives
\[
                  \operatorname{charpol}(\overline M)
                        =(X-\mu)^n.               \tag{G89.9}
\]
It does not assert \(\overline M=\mu I\).

Consider the \(A\)-linear endomorphism of \(M_n(A)\)
\[
                     X\longmapsto cMX-XM.         \tag{G89.10}
\]
After extending the residue field to an algebraic closure,
triangularize \(\overline M\).
The eigenvalues of the induced Sylvester operator are
\(c\mu-\mu=\mu(c-1)\), each with multiplicity \(n^2\).
For example this follows by triangularizing its two commuting
left and right multiplication operators, or by the usual
tensor description of \(\operatorname{End}(k'^n)\).
Its determinant is therefore
\((\mu(c-1))^{n^2}\ne0\).
The determinant of (G89.10) is a unit in \(A\), so the operator
is invertible over \(A\), and also over its fraction field.
An \(X\) satisfying the \(\mathrm{N4}\) equation for every
element of \(D_0\) lies in this one operator's kernel and is zero.
\end{proof}

\begin{theorem}[AC-G89-support: a proper native ordinary determinant quotient]
\label{thm:G89-support}
Retain the actual host, representations, and exact local conclusions
above.  Retain, as explicit source inputs, the propositions
\(\mathrm{LevEx}\), \(\mathrm{ProModularRingExists}\),
\(\mathrm{LamEx}\), and \(\mathrm{ModU}\).
In particular, \(\mathrm{ModU}\) is the proposition defined in
\path{ArithmeticPointClassicality.lean}; it is not proved here
from hypothesis \(A\), and no Lean build is asserted.

Choose \(w\geq2\) with \(w\equiv e\pmod{268}\), and use the
actual representation on the fixed host supplied by AC-G87.
There are a sufficiently large coefficient DVR \(\mathcal O\),
a native source \(\mathrm{HidaSetting}\) with the auxiliary
place of Lemma~\ref{lem:G89-v0}, a native residual datum \(r\)
with characters (G89.2), a universal ordinary determinant ring
\(B\), and its actual inertia-weight ring \(\Lambda\), such that
the ideal \(J\) supplied by \(\mathrm{ModU}\) is proper.
It contains an actual stable arithmetic point:
\[
 J\subset\ker x,\qquad
 x:B\longrightarrow\overline{\mathbb Q}_p,\qquad
 x.D=D_{\rho_\sigma}.                              \tag{G89.11}
\]
Here \(\rho_\sigma\) is in the source normalization, so that
\(\epsilon\rho_\sigma\) is the unitary normalization of AC-G87.

The quotient \(B/J\) is nonzero, finite and torsionfree over
\[
 \Lambda\simeq
       \mathcal O[[X_1,\ldots,X_{3[H^+:\mathbb Q]}]].       \tag{G89.12}
\]
Every minimal prime of \(B/J\) has dimension
\(1+3[H^+:\mathbb Q]\).
The quotient satisfies the literal determinant
unramifiedness conclusion at \(I_0\), and every source point
of \(r\) satisfies \(\mathrm{N4}\).
The source \(\mathrm{P1core}\) and \(\mathrm{P3core}\)
are retained for this \(J\), exactly as supplied by
\(\mathrm{ModU}\).

These conclusions concern the arbitrary polarized cubic residual
datum.  They do not construct the old single-character
\(\mathrm{TwLink}\), an \(\mathrm{NRep}\) of that old frame,
or containment of \(J\) in a prescribed characteristic-\(p\)
curve kernel.
\end{theorem}

\begin{proof}
\emph{The field and coefficient data.}
Choose the excluded finite set in Lemma~\ref{lem:G89-v0}
to contain all ramification of \(H\), of the actual automorphic
representation, and of the fixed residual data.
Every ramified automorphic place is split over \(H^+\)
by AC-G83, and every place above \(p\) is split over \(H^+\).
The extension \(H/H^+\) is unramified at all finite places.
Take \(\mathrm{St}\) to contain the finite set of ramified
places away from \(p\), enlarging it by residual ramification
if necessary; all these places are relatively split.
Together with the completely split residue-degree-one place \(v_0\),
this gives every displayed field and place condition of the
native \(\mathrm{HidaSetting}\).
Apply \(\mathrm{LevEx}\) with its exact set \(S'=S\).

The maximal abelian subextension of \(H/\mathbb Q\) is
multiquadratic, by the commutator calculation above.
Thus \(H\) does not contain \(\mu_{269}\), since its cyclotomic
Galois group has order \(268\) and a quotient of order \(67\).
Choose a finite coefficient field containing the matrices of
\(\rho_\sigma\), its finitely many local flag characters,
the residual characters, and every \(p\)-adic embedding of \(H\).
Its integer ring \(\mathcal O\) satisfies the source
\(\mathrm{Large}\) condition.
This choice is for the present automorphic support construction;
it does not change an earlier deformation coefficient hypothesis.

\emph{The residual datum and its actual integral pair.}
The representation \(\rho_\sigma\) is continuous and has compact
image, so it preserves an \(\mathcal O\)-lattice after the finite
coefficient choice.
Its determinant law and its coefficients are consequently integral.
The raw parameter is
\(\rho_\sigma=\epsilon^{-1}R_w\), with \(R_w\) the unitary
representation of AC-G87.
Its residual determinant is (G89.2), since
\(w\equiv e\pmod{268}\).
The three global characters are distinct: equality of the
endpoints would give \(\overline\gamma^2=1\), while equality
of either endpoint with the middle, followed by CM conjugation,
would give \(\overline\chi^{\,2j}=1\).
These possibilities are excluded by the retained separation.
The conjugate-dual polarization of \(\rho_\sigma\) has
multiplier \(\epsilon^{-2}\), exactly the \(n=3\) multiplier
in the source \(\mathrm{star}\).
It gives the required polarized residual determinant.

At a coefficient-prime place the exact ordered flag has characters
\[
 \kappa_{v,1}=u_{v,1}\chi^j\tau\,\epsilon^{w-1},
 \qquad
 \kappa_{v,2}=u_{v,2}\epsilon^{-1},
 \qquad
 \kappa_{v,3}=u_{v,3}\chi^j\tau\,\epsilon^{-w-1},
                                                        \tag{G89.13}
\]
where the \(u_{v,i}\) are unramified.
Every \(H_v/\mathbb Q_p\) is unramified.
On inertia, the cubic factor has order \(3\), the factor
\(\tau\) has order at most \(2\), and the residual cyclotomic
factor has order \(268\).
The inequality \(268\nmid2e\) therefore distinguishes the
two endpoints, and the cubic factor distinguishes each from
the middle.
Reduction of an integral local flag has the same semisimplified
characters as the restricted residual determinant; these
distinct inertia characters identify their ordering uniquely.
Thus the reductions of (G89.13) are exactly the ordered local
restrictions of (G89.2).

Define \(\mathrm{psib}\) by those restrictions.
Define \(\mathrm{th}\) by restricting the actual \(\kappa_{v,i}\)
to the source group \(\mathrm{mu}\); the values are integral
units and every continuous character kills the closure of the
commutator subgroup, so \(\mathrm{th\_comm}\) holds.
We do not need a coefficient extension as a test object over
a smaller coefficient ring: with this \(\mathcal O\) chosen
from the start, the native test object is \(\mathcal O\)
itself, and \(\mathrm{IsCO.self}\) applies.

The representation factors through the source quotient
\(\mathrm{Lev.pr}\), because it is unramified outside the
chosen conjugation-stable set.
Its determinant \(D\), together with (G89.13), forms an actual
\(\mathrm{Pair}(r,\mathcal O)\).
Here are the individual nonformal checks.
Continuity gives the adic continuity conditions.
The residual determinant and local reductions are those just
verified, and the restrictions to \(\mathrm{mu}\) agree by
the definition of \(\mathrm{th}\).
The polarization gives the source coefficient identities
\(\mathrm{Pair.P}\).
For \(\mathrm{Pair.Ord}\), use the ordered flag:
the rightmost factor
\(\rho(g_3)-\kappa_{v,3}(g_3)\) lowers it by one step,
then the next factor lowers it again, and the final factor
kills it.  Thus the source ordered product \(\mathrm{oEl}\)
is the zero matrix, and its product with any \(\rho(h)\)
has trace zero.
For \(\mathrm{Pair.I}\), the exact AC-G80 types are direct
sums of finite characters of order prime to \(p\), with
zero monodromy, even in repeated-character blocks.
If an inertia element belongs to every residual kernel,
each such character is \(1\): reduction is injective on
roots of unity of order prime to \(p\).
The inertia matrix is therefore the identity.
This is the exact \(\mathrm{Ip}\) condition, not merely
unramifiedness of the residual semisimplification.

The retained \(\mathrm{ProModularRingExists}\) now supplies
\(B\).  Its universal property applied to the pair over
\(\mathcal O\) gives a map \(B\to\mathcal O\);
compose with \(\mathcal O\hookrightarrow\overline{\mathbb Q}_p\)
to obtain \(x\).
This proves the native point and representation identities
required before applying \(\mathrm{PT}\).
Apply \(\mathrm{LamEx}\) to this same \(r\) and \(B\).
It gives precisely (G89.12); the weight map is the source
\(\Lambda.\mathrm{toPair}(B.u)\).

\emph{Applying the actual arithmetic point clause.}
The parallel Hodge--Tate weights of \(\rho_\sigma\), in the
source convention, are
\[
                         (1-w,\,1,\,1+w).         \tag{G89.14}
\]
They are strictly increasing.
The explicit local characters (G89.13) are locally algebraic
with these weights, so \(x\) satisfies the literal
\(\mathrm{Pt.Arith}\) predicate.
The chosen representation is attached to an actual regular
algebraic conjugate-self-dual cuspidal representation, in the
same \(\epsilon^{-1}\) normalization used in
\(\mathrm{Af.IsAutCsd}\).
It is unramified at \(v_0\) by the choice of the excluded set.
Its local flags are the actual flags (G89.13).
Finally no \(\kappa_{v,i}\) equals
\(\kappa_{v,j}\epsilon\): their Hodge--Tate weight differences
belong to \(\{0,\pm w,\pm2w\}\), none of which is \(1\) or
\(-1\) when \(w\ge2\).
This also treats \(i=j\).
All hypotheses of the native \(\mathrm{PT}\) statement
are now verified.

Apply the retained \(\mathrm{ModU}\) input to obtain
an ideal \(J\) with \(\mathrm{Modular}\).
Its \(\mathrm{HidaProps.pt}\) field gives (G89.11).
In particular \(J\ne B\), because \(x(1)=1\).
The determinant at \(x\) is that of an absolutely irreducible
representation.  It is thus an actual stable arithmetic point;
no claim about an \(\mathrm{NRep}\) of a pre-existing
single-character frame is used.

\emph{The remaining ring consequences.}
The fields \(\mathrm{RingProps.fin}\) and
\(\mathrm{RingProps.tf}\) give finite generation and
torsionfreeness of the nonzero \(B/J\) over the actual
\(\Lambda\) in (G89.12).
Let \(Q\) be a minimal prime of \(B/J\).
Since the quotient is Noetherian, \(Q\) is associated.
No nonzero element of the domain \(\Lambda\) is a
zero-divisor on \(B/J\), so \(Q\cap\Lambda=0\).
The finite integral inclusion
\(\Lambda\hookrightarrow(B/J)/Q\) preserves dimension,
which proves the claimed dimension of every minimal component.
The properness already proved is necessary here; it is not
inferred from a vacuous finite quotient.

Lemma~\ref{lem:G89-v0} and
\(\mathrm{RingProps.unr0}\) give the determinant-kernel
unramifiedness at \(I_0\).
Lemma~\ref{lem:G89-N4} gives \(\mathrm{N4}\) for every
source point, not just the selected arithmetic point.
The \(\mathrm{p1}\) and \(\mathrm{p3}\) fields of
\(\mathrm{Modular}\) supply exactly their printed
\(\mathrm{P1core}\) and \(\mathrm{P3core}\) implications.
\end{proof}

\paragraph{Exact source scope and the remaining matching condition.}
The complete literal files used here are
\path{DirichletPrimeToLPartMulLPowerPart.lean},
\path{Counts.lean},
\path{ContinuousResP.lean},
\path{UniversalOrdinaryDeterminantRing.lean},
\path{ArithmeticPointClassicality.lean},
\path{ResMapBijective.lean},
\path{SecondThirdPrimeConditionTransferAlt.lean}, and
\path{SeedPtOf.lean}, at the pinned source revision in the
source manifest.
Their displayed definitions and applied elementary arguments
were read; the retained source propositions are not replaced by
an asserted Lean build or an unproved consequence of \(A\).
The passage from a concrete integral representation to its
\(\mathrm{Pair}\) above is proved directly, so it does not use
the single-character factorization in \path{SeedPtOf.lean}.

The old \(\mathrm{Setup}\) additionally requires
\(\mathrm{TwLink}\), and its
\(\mathrm{IdealSpecializationPoint}\) requires an
\(\mathrm{NRep}\) of the specified frame.
The cubic tuple does not satisfy the unchanged three-power
\(\mathrm{TwLink.chib}\) relation.
Consequently the present proper quotient and its scalar
auxiliary place supply genuine native support data, while the
two-root frame construction, its cohomological inequalities,
and prescribed characteristic-\(p\) curve containment remain
separate requirements.
No quantitative bound for \(v_0^{\mathbb Q}\) is asserted.

\paragraph{Geometric dictionary.}
Continuous finite-dimensional representations unramified outside
the chosen finite set give the corresponding lisse local systems
on the arithmetic open, in the already-audited
SGA~1 V.4/V.7 and XI.6 framework.
Their characteristic polynomials supply the determinant data,
and the ordered local flag supplies the ordinary product
relations.
Here one actual characteristic-zero point proves that the
determinant support is nonempty; it does not prove that an
independently prescribed curve lies in that support.
This paragraph invokes the existing audited dictionary and
claims no new archival reading.
